
\documentclass[12pt]{article}

% =========================
% Geometry & core packages
% =========================
\usepackage[a4paper,margin=0.8in]{geometry}
\usepackage{amsmath,amssymb,amsthm,mathtools}
\usepackage{bm}
\usepackage{array,booktabs}
\usepackage{enumitem}
\usepackage{microtype}
\usepackage{xcolor}
\usepackage{hyperref}
\usepackage{fancyhdr}
\usepackage{draftwatermark}

% =========================
% Watermark
% =========================
\SetWatermarkText{Unofficial -- QRS@NTU -- qrsntu.org}
\SetWatermarkScale{1}
\SetWatermarkLightness{0.99}

% =========================
% Course metadata
% =========================
\newcommand{\CourseCode}{MH1201}
\newcommand{\CourseName}{Linear Algebra II}
\newcommand{\DocType}{Final Exam (Solutions)}
\newcommand{\ExamMeta}{AY 2022/2023, Semester 2}
\newcommand{\ExamMetaShort}{Final AY22-23}

\title{
\Huge \CourseCode\ \CourseName \\[0.3em]
\Large \DocType  \\[0.3em]
\large \ExamMeta
}
\author{
  \textit{\href{https://qrsntu.org}{Quantitative Research Society @NTU}}
}
\date{\today}

% =========================
% Header/footer styles
% =========================
\fancypagestyle{main}{
  \fancyhf{}
  \fancyhead[L]{\CourseCode\ \CourseName\ --\ \ExamMetaShort}
  \fancyhead[R]{\href{https://qrsntu.org}{QRS@NTU}}
  \fancyfoot[C]{\thepage}
  \renewcommand{\headrulewidth}{0.4pt}
  \renewcommand{\footrulewidth}{0pt}
}

\fancypagestyle{plainfirst}{
  \fancyhf{}
  \renewcommand{\headrulewidth}{0pt}
  \renewcommand{\footrulewidth}{0pt}
  \fancyfoot[C]{\scriptsize\textcolor{gray}{\textit{For academic reference only. This document is provided for educational purposes and does not constitute an official question paper, answer key, or any formally authorised academic material. Its content is not guaranteed to be complete or accurate.}}}
}

% =========================
% Convenience macros
% =========================
\newcommand{\R}{\mathbb{R}}
\newcommand{\N}{\mathbb{N}}
\newcommand{\M}{\mathcal{M}}
\newcommand{\V}{\mathcal{V}}
\newcommand{\W}{\mathcal{W}}
\newcommand{\U}{\mathcal{U}}
\newcommand{\B}{\mathcal{B}}
\newcommand{\Ppoly}{\mathcal{P}}
\newcommand{\Null}{\operatorname{null}}
\newcommand{\Range}{\operatorname{range}}
\newcommand{\Span}{\operatorname{span}}
\newcommand{\rank}{\operatorname{rank}}
\newcommand{\diag}{\operatorname{diag}}
\newcommand{\tr}{\operatorname{tr}}
\newcommand{\ip}[2]{\left\langle #1,#2\right\rangle}
\allowdisplaybreaks

% =========================
% Overview block
% =========================
\newenvironment{overview}{
  \par\bigskip
  \begin{center}
    \rule{0.92\linewidth}{0.6pt}\\[-0.2em]
    \textbf{Overview}\\[-0.2em]
    \rule{0.92\linewidth}{0.6pt}
  \end{center}
  \vspace{-0.3em}
  \begin{itemize}[leftmargin=1.6em, itemsep=0.2em, topsep=0.2em]
}{
  \end{itemize}
  \par\bigskip
}

\setlist[enumerate]{itemsep=0.3em, topsep=0.3em}
\setlist[itemize]{itemsep=0.2em, topsep=0.2em}
\setcounter{secnumdepth}{-1}
\setcounter{tocdepth}{3}

\begin{document}

\maketitle
\thispagestyle{plainfirst}
\pagestyle{main}


\begin{overview}
  \item Topics: polynomial maps and parity decomposition, diagonalisation of block matrices, Gram--Schmidt and QR decomposition in $\R^4$, rank--nullity, and recursive block-matrix eigenvectors.
  \item The solutions prioritise exam-efficient computation while still explaining the key linear algebra ideas.
  \item Mark schemes are unofficial and are intended as revision-oriented grading guides.
\end{overview}


\newpage
\tableofcontents
\newpage

% ============================================================
\section{Question 1 \hfill (20 Marks)}

\subsection{Problem}
Let $d$ be a nonnegative integer. Consider the function
\[
T:\Ppoly_d(\R)\to\Ppoly_d(\R),
\qquad
T(f(x))=f(x)+f(-x).
\]
\begin{enumerate}[label=(\alph*)]
  \item Show that $T$ is a linear map.
  \item Fix $d=2$.
  \begin{enumerate}[label=(\roman*)]
    \item Determine the standard matrix for $T$.
    \item Determine a basis for $\Null(T)$.
    \item Determine a basis for $\Range(T)$.
  \end{enumerate}
  \item Show that
  \[
  \dim\Null(T)=\left\lceil\frac d2\right\rceil,
  \qquad
  \dim\Range(T)=\left\lfloor\frac d2\right\rfloor+1
  \]
  for all $d$.
\end{enumerate}

\subsection{Solution}

\subsubsection{Method 1: Separate even and odd powers}
For $f,g\in\Ppoly_d(\R)$ and $c\in\R$,
\[
T(f+g)(x)=(f+g)(x)+(f+g)(-x)=T(f)(x)+T(g)(x),
\]
and
\[
T(cf)(x)=cf(x)+cf(-x)=cT(f)(x).
\]
Hence $T$ is linear.

For $d=2$, use the ordered standard basis $(1,x,x^2)$. If
\[
f(x)=a_0+a_1x+a_2x^2,
\]
then
\[
f(-x)=a_0-a_1x+a_2x^2,
\]
so
\[
T(f)(x)=2a_0+2a_2x^2.
\]
Therefore
\[
\boxed{[T]=\begin{bmatrix}2&0&0\\0&0&0\\0&0&2\end{bmatrix}.}
\]
The nullspace consists of polynomials whose even part is zero, so
\[
\boxed{\Null(T)=\Span\{x\}.}
\]
The range consists of even polynomials of degree at most $2$, so
\[
\boxed{\Range(T)=\Span\{1,x^2\}.}
\]

For general $d$, write
\[
f(x)=a_0+a_1x+\cdots+a_dx^d.
\]
Then
\[
T(f)(x)=2(a_0+a_2x^2+a_4x^4+\cdots).
\]
Thus $\Null(T)$ consists exactly of odd polynomials of degree at most $d$, and $\Range(T)$ consists exactly of even polynomials of degree at most $d$.

The odd exponents among $0,1,\ldots,d$ are
\[
1,3,5,\ldots\le d,
\]
whose count is $\lceil d/2\rceil$. The even exponents are
\[
0,2,4,\ldots\le d,
\]
whose count is $\lfloor d/2\rfloor+1$. Hence
\[
\boxed{\dim\Null(T)=\left\lceil\frac d2\right\rceil,
\qquad
\dim\Range(T)=\left\lfloor\frac d2\right\rfloor+1.}
\]

\subsubsection{Method 2: Projection onto the even subspace}
Every polynomial $f\in\Ppoly_d(\R)$ can be written uniquely as
\[
f=f_{\mathrm{even}}+f_{\mathrm{odd}},
\]
where
\[
f_{\mathrm{even}}(x)=\frac{f(x)+f(-x)}2,
\qquad
f_{\mathrm{odd}}(x)=\frac{f(x)-f(-x)}2.
\]
Then
\[
T(f)=2f_{\mathrm{even}}.
\]
So $T$ kills the odd subspace and maps onto the even subspace. Counting basis monomials gives the same dimension formulae.

\subsection{Mark Scheme (20 marks)}
\begin{itemize}[leftmargin=1.6em]
  \item \textbf{(a) Linearity (4 marks)}
  \begin{itemize}[leftmargin=1.6em]
    \item (2) Correct additivity proof.
    \item (1) Correct homogeneity proof.
    \item (1) States that $T$ is linear.
  \end{itemize}
  \item \textbf{(b) Case $d=2$ (9 marks)}
  \begin{itemize}[leftmargin=1.6em]
    \item (3) Correct standard matrix relative to $(1,x,x^2)$.
    \item (3) Correct basis for $\Null(T)$.
    \item (3) Correct basis for $\Range(T)$.
  \end{itemize}
  \item \textbf{(c) General dimensions (7 marks)}
  \begin{itemize}[leftmargin=1.6em]
    \item (3) Correctly identifies the nullspace as the odd-polynomial subspace.
    \item (3) Correctly identifies the range as the even-polynomial subspace.
    \item (1) Correctly counts and states both dimension formulae.
  \end{itemize}
\end{itemize}

% ============================================================
\newpage
\section{Question 2 \hfill (20 Marks)}

\subsection{Problem}
Let $0<a\le1$ and $0<b\le1$, where $a$ and $b$ are real numbers. Consider
\[
A=\begin{bmatrix}
2&0&0\\
0&a&1-b\\
0&1-a&b
\end{bmatrix}.
\]
\begin{enumerate}[label=(\alph*)]
  \item Show that the characteristic polynomial for $A$ is $(2-\lambda)(1-\lambda)(a+b-1-\lambda)$. Hence determine the eigenvalues of $A$ in terms of $a$ and $b$.
  \item Show that $A$ is diagonalizable.
  \item Fix $a=b$ and set $a<1$.
  \begin{enumerate}[label=(\roman*)]
    \item Find a diagonal matrix $D$ and an invertible matrix $P$ such that $A=PDP^{-1}$.
    \item Fix $0.25<a<0.75$. Approximate all entries in $A^{20}$ to the nearest one decimal place.
  \end{enumerate}
\end{enumerate}

\subsection{Solution}

\subsubsection{Method 1: Block diagonalisation}
Since $A$ is block diagonal with a $1\times1$ block $[2]$ and a lower block
\[
B=\begin{bmatrix}a&1-b\\1-a&b\end{bmatrix},
\]
we have
\[
\det(A-\lambda I)=(2-\lambda)\det\begin{bmatrix}a-\lambda&1-b\\1-a&b-\lambda\end{bmatrix}.
\]
Now
\begin{align*}
\det\begin{bmatrix}a-\lambda&1-b\\1-a&b-\lambda\end{bmatrix}
&=(a-\lambda)(b-\lambda)-(1-a)(1-b)\\
&=\lambda^2-(a+b)\lambda+(a+b-1)\\
&=(1-\lambda)(a+b-1-\lambda).
\end{align*}
Thus
\[
\boxed{\det(A-\lambda I)=(2-\lambda)(1-\lambda)(a+b-1-\lambda)}
\]
and the eigenvalues are
\[
\boxed{2,
\quad 1,
\quad a+b-1.}
\]

The eigenvalue $2$ is separated from the lower block. The lower block $B$ has eigenvalues $1$ and $a+b-1$. If they are distinct, then $B$ is diagonalizable. If they are equal, then
\[
1=a+b-1 \quad\Longrightarrow\quad a+b=2.
\]
Since $0<a\le1$ and $0<b\le1$, this forces $a=b=1$, in which case
\[
B=I_2,
\]
which is diagonalizable. Hence $B$ is always diagonalizable, and so
\[
\boxed{A\text{ is diagonalizable}.}
\]

Now fix $a=b<1$. Then
\[
A=\begin{bmatrix}
2&0&0\\
0&a&1-a\\
0&1-a&a
\end{bmatrix}.
\]
The lower block has eigenvectors $(1,1)^T$ and $(1,-1)^T$ with eigenvalues $1$ and $2a-1$, respectively. Therefore take
\[
v_1=\begin{bmatrix}1\\0\\0\end{bmatrix},
\quad
v_2=\begin{bmatrix}0\\1\\1\end{bmatrix},
\quad
v_3=\begin{bmatrix}0\\1\\-1\end{bmatrix}.
\]
Thus
\[
\boxed{
P=\begin{bmatrix}1&0&0\\0&1&1\\0&1&-1\end{bmatrix},
\qquad
D=\diag(2,1,2a-1),
\qquad A=PDP^{-1}.}
\]
Then
\[
A^{20}=PD^{20}P^{-1}.
\]
Let
\[
r=(2a-1)^{20}.
\]
The lower $2\times2$ block of $A^{20}$ is
\[
\frac12\begin{bmatrix}1+r&1-r\\1-r&1+r\end{bmatrix}.
\]
Hence
\[
A^{20}=\begin{bmatrix}
2^{20}&0&0\\
0&\frac{1+r}{2}&\frac{1-r}{2}\\
0&\frac{1-r}{2}&\frac{1+r}{2}
\end{bmatrix}.
\]
If $0.25<a<0.75$, then $|2a-1|<0.5$, so
\[
|r|<0.5^{20}\approx9.54\times10^{-7}.
\]
Thus $r$ is $0.0$ to one decimal place. Since $2^{20}=1{,}048{,}576$,
\[
\boxed{
A^{20}\approx
\begin{bmatrix}
1{,}048{,}576.0&0.0&0.0\\
0.0&0.5&0.5\\
0.0&0.5&0.5
\end{bmatrix}.}
\]

\subsubsection{Method 2: Use symmetry of the lower block when $a=b$}
For part (c), when $a=b$, the lower block is
\[
\begin{bmatrix}a&1-a\\1-a&a\end{bmatrix}
=a\begin{bmatrix}1&0\\0&1\end{bmatrix}+(1-a)\begin{bmatrix}0&1\\1&0\end{bmatrix}.
\]
The vectors $(1,1)^T$ and $(1,-1)^T$ diagonalize the swap matrix. This gives the same $P$ and $D$ immediately, and is the fastest way to compute $A^{20}$ in part (c).

\subsection{Mark Scheme (20 marks)}
\begin{itemize}[leftmargin=1.6em]
  \item \textbf{(a) Characteristic polynomial and eigenvalues (5 marks)}
  \begin{itemize}[leftmargin=1.6em]
    \item (2) Correct block determinant setup.
    \item (2) Correct expansion and factorisation.
    \item (1) Correct eigenvalues.
  \end{itemize}
  \item \textbf{(b) Diagonalizability (5 marks)}
  \begin{itemize}[leftmargin=1.6em]
    \item (2) Analyses the lower $2\times2$ block and its eigenvalues.
    \item (2) Handles the repeated case $a=b=1$ correctly.
    \item (1) Concludes that $A$ is diagonalizable under the stated conditions.
  \end{itemize}
  \item \textbf{(c) Diagonalisation and power (10 marks)}
  \begin{itemize}[leftmargin=1.6em]
    \item (3) Finds correct eigenvectors for $a=b<1$.
    \item (2) Gives valid $P$ and $D$.
    \item (2) Derives the formula for $A^{20}$ using $PD^{20}P^{-1}$.
    \item (2) Correctly uses $|2a-1|<0.5$ to approximate $r$ as $0.0$ to one decimal place.
    \item (1) Gives all nine entries of the approximated matrix.
  \end{itemize}
\end{itemize}

% ============================================================
\newpage
\section{Question 3 \hfill (20 Marks)}

\subsection{Problem}
\begin{enumerate}[label=(\alph*)]
  \item Consider the following vectors in $\R^4$ with the usual Euclidean inner product:
  \[
  u_1=\begin{bmatrix}1\\-1\\0\\0\end{bmatrix},
  \qquad
  u_2=\begin{bmatrix}0\\1\\-1\\0\end{bmatrix},
  \qquad
  u_3=\begin{bmatrix}0\\0\\1\\-1\end{bmatrix}.
  \]
  Apply the Gram--Schmidt procedure to $\{u_1,u_2,u_3\}$ to obtain an orthogonal basis for $\Span\{u_1,u_2,u_3\}$.
  \item Let
  \[
  A=\begin{bmatrix}
  1&0&0&1\\
  -1&1&0&1\\
  0&-1&1&1\\
  0&0&-1&1
  \end{bmatrix}.
  \]
  Find the QR-decomposition for $A$.
\end{enumerate}

\subsection{Solution}

\subsubsection{Method 1: Gram--Schmidt column by column}

For part (a), we only need an orthogonal basis, not an orthonormal basis.

First,
\[
v_1=u_1=
\begin{bmatrix}1\\-1\\0\\0\end{bmatrix}.
\]

Next,
\[
v_2
=
u_2-\frac{u_2\cdot v_1}{v_1\cdot v_1}v_1.
\]
Since
\[
u_2\cdot v_1=-1,
\qquad
v_1\cdot v_1=2,
\]
we get
\[
v_2
=
\begin{bmatrix}0\\1\\-1\\0\end{bmatrix}
-\left(-\frac12\right)
\begin{bmatrix}1\\-1\\0\\0\end{bmatrix}
=
\begin{bmatrix}1/2\\1/2\\-1\\0\end{bmatrix}.
\]
Multiplying by $2$ gives the simpler orthogonal vector
\[
v_2=
\begin{bmatrix}1\\1\\-2\\0\end{bmatrix}.
\]

Finally,
\[
v_3
=
u_3
-\frac{u_3\cdot v_1}{v_1\cdot v_1}v_1
-\frac{u_3\cdot v_2}{v_2\cdot v_2}v_2.
\]
Now
\[
u_3\cdot v_1=0,
\qquad
u_3\cdot v_2=-2,
\qquad
v_2\cdot v_2=6.
\]
Hence
\[
v_3
=
\begin{bmatrix}0\\0\\1\\-1\end{bmatrix}
-\left(-\frac13\right)
\begin{bmatrix}1\\1\\-2\\0\end{bmatrix}
=
\begin{bmatrix}1/3\\1/3\\1/3\\-1\end{bmatrix}.
\]
Multiplying by $3$ gives
\[
v_3=
\begin{bmatrix}1\\1\\1\\-3\end{bmatrix}.
\]

Therefore an orthogonal basis for $\Span\{u_1,u_2,u_3\}$ is
\[
\boxed{
\left\{
\begin{bmatrix}1\\-1\\0\\0\end{bmatrix},
\begin{bmatrix}1\\1\\-2\\0\end{bmatrix},
\begin{bmatrix}1\\1\\1\\-3\end{bmatrix}
\right\}.}
\]

For part (b), we need an orthonormal basis in order to form the matrix $Q$.

Normalising the three orthogonal vectors above gives
\[
e_1=\frac1{\sqrt2}
\begin{bmatrix}1\\-1\\0\\0\end{bmatrix},
\qquad
e_2=\frac1{\sqrt6}
\begin{bmatrix}1\\1\\-2\\0\end{bmatrix},
\qquad
e_3=\frac1{2\sqrt3}
\begin{bmatrix}1\\1\\1\\-3\end{bmatrix}.
\]

The first three columns of $A$ are $u_1,u_2,u_3$. The fourth column is
\[
u_4=\begin{bmatrix}1\\1\\1\\1\end{bmatrix}.
\]
This vector is orthogonal to $u_1,u_2,u_3$, and hence also orthogonal to $e_1,e_2,e_3$. Therefore
\[
e_4=\frac12
\begin{bmatrix}1\\1\\1\\1\end{bmatrix}.
\]

Taking
\[
Q=[e_1\ e_2\ e_3\ e_4],
\]
we get
\[
Q=
\begin{bmatrix}
\frac1{\sqrt2}&\frac1{\sqrt6}&\frac1{2\sqrt3}&\frac12\\[0.3em]
-\frac1{\sqrt2}&\frac1{\sqrt6}&\frac1{2\sqrt3}&\frac12\\[0.3em]
0&-\frac2{\sqrt6}&\frac1{2\sqrt3}&\frac12\\[0.3em]
0&0&-\frac3{2\sqrt3}&\frac12
\end{bmatrix}.
\]

Then
\[
R=Q^TA.
\]
Computing the entries gives
\[
R=
\begin{bmatrix}
\sqrt2&-\frac1{\sqrt2}&0&0\\[0.3em]
0&\frac3{\sqrt6}&-\frac2{\sqrt6}&0\\[0.3em]
0&0&\frac2{\sqrt3}&0\\[0.3em]
0&0&0&2
\end{bmatrix}.
\]

Therefore
\[
\boxed{A=QR.}
\]

\subsubsection{Method 2: Orthogonal complement shortcut}

From part (a), an orthogonal basis for $\Span\{u_1,u_2,u_3\}$ is
\[
\left\{
\begin{bmatrix}1\\-1\\0\\0\end{bmatrix},
\begin{bmatrix}1\\1\\-2\\0\end{bmatrix},
\begin{bmatrix}1\\1\\1\\-3\end{bmatrix}
\right\}.
\]
For the QR decomposition, these vectors are then normalised to form the first three columns of $Q$.

The fourth column of $A$ is
\[
u_4=
\begin{bmatrix}1\\1\\1\\1\end{bmatrix},
\]
which is visibly orthogonal to $u_1,u_2,u_3$. Hence it is also orthogonal to the first three Gram--Schmidt vectors. Normalising it immediately gives
\[
e_4=\frac12
\begin{bmatrix}1\\1\\1\\1\end{bmatrix}.
\]
Then the entries of $R$ are obtained from
\[
R_{ij}=e_i\cdot u_j.
\]
This avoids solving any additional linear equations for the fourth column.

\subsection{Mark Scheme (20 marks)}
\begin{itemize}[leftmargin=1.6em]
  \item \textbf{(a) Gram--Schmidt for an orthogonal basis (10 marks)}
  \begin{itemize}[leftmargin=1.6em]
    \item (2) Correctly identifies $v_1=u_1$.
    \item (3) Correctly computes the second orthogonal vector $v_2$.
    \item (3) Correctly computes the third orthogonal vector $v_3$.
    \item (2) Presents the final orthogonal basis clearly, without unnecessarily normalising it.
  \end{itemize}

  \item \textbf{(b) QR decomposition (10 marks)}
  \begin{itemize}[leftmargin=1.6em]
    \item (2) Correctly normalises the orthogonal basis vectors from part (a).
    \item (2) Identifies and normalises the fourth orthogonal column.
    \item (2) Forms the correct orthogonal matrix $Q$.
    \item (3) Computes the correct upper-triangular matrix $R$.
    \item (1) States $A=QR$ with correct dimensions and triangular form.
  \end{itemize}
\end{itemize}

% ============================================================
\newpage
\section{Question 4 \hfill (20 Marks)}

\subsection{Problem}
Let $\V$ be a finite-dimensional vector space and consider a linear map $T:\V\to\R$. Suppose $\{u_1,u_2,\ldots,u_s\}$ is a basis for $\Null(T)$. Show that if $T(v)\neq0$, then
\[
\{u_1,u_2,\ldots,u_s,v\}
\]
is a basis for $\V$.

\subsection{Solution}

\subsubsection{Method 1: Rank--nullity plus linear independence}
Since $T(v)\neq0$, $T$ is not the zero map. Because the codomain is $\R$, the range is a nonzero subspace of $\R$, so
\[
\dim\Range(T)=1.
\]
By rank--nullity,
\[
\dim\V=\dim\Null(T)+\dim\Range(T)=s+1.
\]
Now suppose
\[
a_1u_1+\cdots+a_su_s+bv=0.
\]
Apply $T$ to both sides:
\[
0=a_1T(u_1)+\cdots+a_sT(u_s)+bT(v)=bT(v).
\]
Since $T(v)\neq0$, we get $b=0$. Then
\[
a_1u_1+\cdots+a_su_s=0.
\]
Because $u_1,
\ldots,u_s$ are a basis for $\Null(T)$, they are linearly independent, so all $a_i=0$. Thus the set is linearly independent. It has $s+1=\dim\V$ vectors, so it is a basis for $\V$.

\[
\boxed{\{u_1,u_2,\ldots,u_s,v\}\text{ is a basis for }\V.}
\]

\subsubsection{Method 2: Spanning by reducing to the kernel}
For any $x\in\V$, define
\[
y=x-\frac{T(x)}{T(v)}v.
\]
Then
\[
T(y)=T(x)-\frac{T(x)}{T(v)}T(v)=0,
\]
so $y\in\Null(T)$. Since $u_1,\ldots,u_s$ span $\Null(T)$,
\[
y=a_1u_1+\cdots+a_su_s.
\]
Thus
\[
x=a_1u_1+\cdots+a_su_s+\frac{T(x)}{T(v)}v,
\]
so $u_1,\ldots,u_s,v$ span $\V$. The same linear-independence argument from Method 1 shows the set is independent. Hence it is a basis.

\subsection{Mark Scheme (20 marks)}
\begin{itemize}[leftmargin=1.6em]
  \item (3) Correctly observes that $T(v)\neq0$ implies $\Range(T)=\R$ and hence $\dim\Range(T)=1$.
  \item (3) Applies rank--nullity to obtain $\dim\V=s+1$.
  \item (4) Sets up a general linear relation among $u_1,\ldots,u_s,v$.
  \item (4) Applies $T$ to show the coefficient of $v$ is zero.
  \item (3) Uses independence of the kernel basis to show all remaining coefficients are zero.
  \item (2) Concludes that a linearly independent set with $\dim\V$ vectors is a basis.
  \item (1) Presents the final conclusion clearly.
\end{itemize}

% ============================================================
\newpage
\section{Question 5 \hfill (20 Marks)}

\subsection{Problem}
Let
\[
H^{(1)}=\begin{bmatrix}2&-3\\1&-2\end{bmatrix}.
\]
\begin{enumerate}[label=(\alph*)]
  \item Let $e_1^{(1)}=(1,1)^T$ and $e_2^{(1)}=(3,1)^T$. Explain why $e_1^{(1)}$ and $e_2^{(1)}$ are eigenvectors of $H^{(1)}$, and determine their eigenvalues.
  \item Define
  \[
  H^{(2)}=\begin{bmatrix}2H^{(1)}&-3H^{(1)}\\H^{(1)}&-2H^{(1)}\end{bmatrix}.
  \]
  Set
  \[
  e_{11}^{(2)}=\begin{bmatrix}e_1^{(1)}\\e_1^{(1)}\end{bmatrix},
  \quad
  e_{12}^{(2)}=\begin{bmatrix}e_2^{(1)}\\e_2^{(1)}\end{bmatrix},
  \quad
  e_{21}^{(2)}=\begin{bmatrix}3e_1^{(1)}\\e_1^{(1)}\end{bmatrix},
  \quad
  e_{22}^{(2)}=\begin{bmatrix}3e_2^{(1)}\\e_2^{(1)}\end{bmatrix}.
  \]
  Explain why these four vectors are eigenvectors of $H^{(2)}$ and determine their eigenvalues.
  \item Let $P^{(1)}=[e_1^{(1)}\ e_2^{(1)}]$ and $P^{(2)}=[e_{11}^{(2)}\ e_{12}^{(2)}\ e_{21}^{(2)}\ e_{22}^{(2)}]$. Show that
  \[
  P^{(2)}=\begin{bmatrix}P^{(1)}&0\\0&P^{(1)}\end{bmatrix}
  \begin{bmatrix}I_2&3I_2\\I_2&I_2\end{bmatrix}.
  \]
  \item For $k\ge3$, recursively define
  \[
  H^{(k)}=\begin{bmatrix}2H^{(k-1)}&-3H^{(k-1)}\\H^{(k-1)}&-2H^{(k-1)}\end{bmatrix}.
  \]
  Show that $1$ and $-1$ are eigenvalues of $H^{(k)}$, and the multiplicities for both $1$ and $-1$ are $2^{k-1}$. You may assume that
  \[
  \begin{bmatrix}I_m&3I_m\\I_m&I_m\end{bmatrix}
  \]
  is invertible for every $m\ge1$.
\end{enumerate}

\subsection{Solution}

\subsubsection{Method 1: Eigenvector propagation through block matrices}
First,
\[
H^{(1)}e_1^{(1)}=\begin{bmatrix}2&-3\\1&-2\end{bmatrix}
\begin{bmatrix}1\\1\end{bmatrix}
=\begin{bmatrix}-1\\-1\end{bmatrix}=-e_1^{(1)},
\]
so $e_1^{(1)}$ has eigenvalue $-1$. Also
\[
H^{(1)}e_2^{(1)}=\begin{bmatrix}2&-3\\1&-2\end{bmatrix}
\begin{bmatrix}3\\1\end{bmatrix}
=\begin{bmatrix}3\\1\end{bmatrix}=e_2^{(1)},
\]
so $e_2^{(1)}$ has eigenvalue $1$.

Now let $x$ be an eigenvector of $H^{(1)}$ with eigenvalue $\lambda$, and let $(\alpha,\beta)^T$ be an eigenvector of $H^{(1)}$ with eigenvalue $\mu$. Then
\[
H^{(2)}\begin{bmatrix}\alpha x\\\beta x\end{bmatrix}
=\begin{bmatrix}(2\alpha-3\beta)H^{(1)}x\\(\alpha-2\beta)H^{(1)}x\end{bmatrix}
=\lambda\begin{bmatrix}(2\alpha-3\beta)x\\(\alpha-2\beta)x\end{bmatrix}
=\lambda\mu\begin{bmatrix}\alpha x\\\beta x\end{bmatrix}.
\]
Thus the new eigenvalue is the product $\lambda\mu$.

Applying this to the four listed vectors gives
\[
\boxed{e_{11}^{(2)}\text{ has eigenvalue }1,\quad
 e_{12}^{(2)}\text{ has eigenvalue }-1,\quad
 e_{21}^{(2)}\text{ has eigenvalue }-1,\quad
 e_{22}^{(2)}\text{ has eigenvalue }1.}
\]

For part (c), the first block matrix places $P^{(1)}$ in the top and bottom blocks, while the second block matrix chooses either the coefficient pair $(1,1)^T$ or $(3,1)^T$. Multiplying the two block matrices gives columns
\[
\begin{bmatrix}e_1^{(1)}\\e_1^{(1)}\end{bmatrix},
\begin{bmatrix}e_2^{(1)}\\e_2^{(1)}\end{bmatrix},
\begin{bmatrix}3e_1^{(1)}\\e_1^{(1)}\end{bmatrix},
\begin{bmatrix}3e_2^{(1)}\\e_2^{(1)}\end{bmatrix},
\]
which are precisely the columns of $P^{(2)}$. Hence
\[
\boxed{P^{(2)}=\begin{bmatrix}P^{(1)}&0\\0&P^{(1)}\end{bmatrix}
\begin{bmatrix}I_2&3I_2\\I_2&I_2\end{bmatrix}.}
\]

For general $k$, suppose $x$ is an eigenvector of $H^{(k-1)}$ with eigenvalue $\lambda\in\{1,-1\}$. Then
\[
H^{(k)}\begin{bmatrix}x\\x\end{bmatrix}
=\begin{bmatrix}-H^{(k-1)}x\\-H^{(k-1)}x\end{bmatrix}
=-\lambda\begin{bmatrix}x\\x\end{bmatrix},
\]
and
\[
H^{(k)}\begin{bmatrix}3x\\x\end{bmatrix}
=\begin{bmatrix}3H^{(k-1)}x\\H^{(k-1)}x\end{bmatrix}
=\lambda\begin{bmatrix}3x\\x\end{bmatrix}.
\]
Thus each eigenvector at level $k-1$ produces one eigenvector with the same eigenvalue and one with the opposite eigenvalue at level $k$.

Let $P^{(k-1)}$ be an eigenvector matrix for $H^{(k-1)}$ and put $m=2^{k-1}$. Define
\[
P^{(k)}=\begin{bmatrix}P^{(k-1)}&0\\0&P^{(k-1)}\end{bmatrix}
\begin{bmatrix}I_m&3I_m\\I_m&I_m\end{bmatrix}.
\]
By induction $P^{(k-1)}$ is invertible, and the second block matrix is assumed invertible. Therefore $P^{(k)}$ is invertible, so these constructed eigenvectors form a basis of $\R^{2^k}$.

At level $k-1$, there are $2^{k-2}$ eigenvectors with eigenvalue $1$ and $2^{k-2}$ with eigenvalue $-1$. Each group produces one eigenvector for eigenvalue $1$ and one for eigenvalue $-1$ at level $k$. Therefore
\[
\boxed{\text{multiplicity of }1=2^{k-1},
\qquad
\text{multiplicity of }-1=2^{k-1}.}
\]

\subsubsection{Method 2: Kronecker-product interpretation}
The recursion may be viewed as
\[
H^{(k)}=H^{(1)}\otimes H^{(k-1)}
\]
up to the same block convention. Eigenvalues of such block products multiply. Since $H^{(1)}$ has eigenvalues $1$ and $-1$, every recursion step splits the existing eigenbasis into one copy with the same sign and one copy with the opposite sign. This gives the same multiplicity doubling argument as Method 1. This method is optional, since the block-vector calculation in Method 1 is usually more exam-safe.

\subsection{Mark Scheme (20 marks)}
\begin{itemize}[leftmargin=1.6em]
  \item \textbf{(a) Base eigenvectors (4 marks)}
  \begin{itemize}[leftmargin=1.6em]
    \item (2) Verifies $e_1^{(1)}$ with eigenvalue $-1$.
    \item (2) Verifies $e_2^{(1)}$ with eigenvalue $1$.
  \end{itemize}
  \item \textbf{(b) Level-two eigenvectors (5 marks)}
  \begin{itemize}[leftmargin=1.6em]
    \item (2) Establishes the block-vector product rule.
    \item (2) Applies it correctly to the four listed vectors.
    \item (1) States all four eigenvalues correctly.
  \end{itemize}
  \item \textbf{(c) Matrix factorisation (4 marks)}
  \begin{itemize}[leftmargin=1.6em]
    \item (2) Correctly identifies how the first block matrix inserts $P^{(1)}$.
    \item (2) Correctly verifies the columns of $P^{(2)}$.
  \end{itemize}
  \item \textbf{(d) General recursion and multiplicities (7 marks)}
  \begin{itemize}[leftmargin=1.6em]
    \item (2) Shows how $[x;x]$ and $[3x;x]$ transform under $H^{(k)}$.
    \item (2) Constructs or describes an eigenbasis using the invertible block matrix.
    \item (2) Correctly counts $2^{k-1}$ eigenvectors for each eigenvalue.
    \item (1) States the final multiplicity conclusion clearly.
  \end{itemize}
\end{itemize}

\end{document}
